\chapter{Background: Mako and the Original Paxos Path}

\Mako{} is a geo-replicated transactional key-value store designed for high
throughput under strong consistency \cite{mako-osdi25}. The system is built for
online transaction processing workloads, where a server must execute many small
transactions per second while still making each transaction durable and
replicated. The central tension is that transaction execution can use many CPU
cores, but replication can easily become a narrow point if all completed work
has to pass through one serial path.

This chapter focuses on the part of \Mako{} that matters for this thesis: the
execute-replicate-replay pipeline. It does not attempt to re-explain all of
\Mako{}'s storage engine, concurrency-control protocol, or geo-replication
policy. The goal is narrower. A reader should leave this chapter understanding
what work the original \Paxos{} backend performed, why its shape was good for a
multicore transaction system, and what must remain true if \Raft{} replaces it.

\section{Transactions, Commit, and Replay}

A transaction is a group of reads and writes that the database treats as one
logical operation. In a simple example, transferring money between two accounts
requires deducting from one account and adding to another. The database should
not expose only half of that operation. Either the whole transaction becomes
part of the database history, or none of it does.

In \Mako{}, transaction work is first submitted to a leader replica. The leader
is the process that accepts the benchmark or client-side transaction work. Its
worker threads execute transactions in parallel: they run the transaction
logic, check the database state needed by the transaction, and produce the
resulting update. The important artifact for replication is not a copy of the
entire database. It is a transaction log record: a compact serialized record
describing the effects that must become visible once the transaction is made
official.

Commit is the point where the replication protocol makes a transaction log
record part of the agreed ordered history. This does not mean every follower has
already replayed the transaction into its local database. In a replicated log
protocol, commit means the log record has been accepted according to the
protocol's rules, such as a quorum agreement in \Paxos{} or \Raft{}. Once a
record is committed, replicas must apply the same committed history in the same
order for the relevant stream.

Replay is the follower-side application of an already committed log record.
The follower reads the committed record and applies its effects to the
follower's local database state. Replay is therefore different from transaction
execution. The leader has already decided what the transaction's effects are.
The follower is not re-running the full transaction as a new speculative
operation; it is applying a committed description of that operation's effects.

This distinction is important for the evaluation later in the thesis. A system
can report high leader-side committed throughput while followers are not keeping
up with replay. That result is not healthy end-to-end replication. For this
reason, the thesis reports both leader-side committed throughput and follower
replay throughput when evaluating the \Raft{} designs.

\section{Worker Lanes and Replication Streams}

\Mako{} inherits the broad systems idea that appears in ROLIS: multicore
transaction execution should not be forced through one global replication lane
when independent worker lanes can make progress in parallel
\cite{rolis-eurosys22}. In this thesis, a partition means the execution and
replication lane associated with a worker thread and its transaction-log
stream. The term is used from the perspective of the benchmark and replication
architecture: it identifies the lane whose committed records must be ordered
consistently.

A worker lane is the path of work produced by one Mako worker thread: the
worker executes transactions and emits a sequence of transaction log records.
A replication stream is the ordered sequence of records for that lane. The
consensus backend decides the order within each stream, so record 1 in a lane
precedes record 2 in that same lane, and replicas replay the committed records
in that order. Different lanes can still progress independently when the
replication backend gives them separate streams. The original Multi-Paxos path
does not require every lane to wait behind one global order; it preserves the
order that matters within each lane while allowing other lanes to advance.

This organization gives the original system its performance shape. Many worker
threads can execute transactions, and the replication layer can preserve that
parallelism by ordering different partition streams independently. The point is
not merely that \Paxos{} is the original protocol. The point is that the
original backend gave \Mako{} multiple ordered streams rather than forcing all
worker-generated records through one stream.

This distinction is the easiest way to understand the Multi-Raft and
Single-Raft designs later in the thesis. Multi-Raft keeps the same idea:
separate worker lanes continue to have separate ordered replication streams.
Single-Raft deliberately changes the shape by routing all lanes through one
consolidated Raft log. That simpler shape can still be correct, but it must be
measured carefully because one shared stream can create pressure in the apply
and replay stages even when leader-side commit throughput looks high.

That shape is the baseline that the \Raft{} implementation must confront. A
direct replacement should ask whether \Raft{} can preserve the same
per-partition parallelism. A more consolidated replacement can simplify the
replication structure, but then it must show that the single ordered path does
not hide another bottleneck in apply or replay.

\section{Original Multi-Paxos Path}

The original \Mako{} implementation used a Multi-Paxos backend to order
transaction log records. The leader-side workers produced log records, the
replication interface submitted those records to the \Paxos{} backend, and
followers replayed committed records into their local database state.

\figref{fig:mako-baseline-replication} shows this path at the level needed for
the rest of the thesis. The left side is the leader process. Its Mako worker
threads execute transactions and produce serialized transaction log records.
Those records are submitted through a common replication interface rather than
being hard-coded directly to one consensus implementation. The center of the
figure is the original Multi-Paxos backend. It establishes the committed log
order for each partition stream. The right side shows follower processes. They
receive committed log records and replay them into their own local database
state.

\begin{figure}[htbp]
\centering
\includegraphics[width=0.96\linewidth]{figures/manual/fig01_mako_baseline_replication.pdf}
\caption[Original Mako replication path]{Original Mako replication path. Leader workers submit transaction log records through the shared replication interface, and follower replicas replay committed logs into local database state.}
\label{fig:mako-baseline-replication}
\end{figure}

The figure also shows why this thesis is not simply swapping one library call
for another. The replication backend sits in the middle of the transaction
pipeline. If it changes the ordering shape or slows replay, the change affects
the end-to-end system. The thesis therefore treats \Raft{} integration as a
systems problem: preserve the useful parts of \Mako{}'s original pipeline while
changing the consensus protocol behind the replication boundary.

\section{The Replication Boundary}

The ability to replace \Paxos{} comes from \Mako{}'s replication boundary. The
transaction execution layer does not need to know whether \Paxos{} or \Raft{}
orders a log record. It submits a record to a helper layer, and the configured
backend handles leader election, log replication, commit advancement, and
delivery of committed records.

This boundary is valuable for two reasons. First, it makes the implementation
tractable: \Raft{} can be integrated as a backend rather than by rewriting
\Mako{}'s transaction execution engine. Second, it makes the evaluation fair.
The benchmark workload, worker-thread counts, CPU accounting, and result
extraction can remain the same while the replication backend changes. When the
evaluation later compares \Paxos{}, Multi-Raft, and Single-Raft, it is comparing
different backend designs under a shared outer transaction workload.

The boundary is not a performance guarantee by itself. A backend can satisfy
the interface and still have the wrong shape for \Mako{}. For example, a
backend that serializes too much work through one apply or replay stage can
look acceptable at the point where the leader commits records, while follower
replay falls behind. This is why the thesis evaluates both the protocol shape
and the replay path.

\section{Modes Evaluated in This Thesis}

\tabref{tab:modes} summarizes the replication modes used throughout the
evaluation. The table is included here because the names appear repeatedly in
the rest of the thesis.

\begin{table}[htbp]
\centering
\caption[Replication modes]{Replication modes evaluated in this thesis.}
\label{tab:modes}
\small
\begin{tabularx}{\linewidth}{@{}>{\raggedright\arraybackslash}p{0.28\linewidth}>{\raggedright\arraybackslash}X@{}}
\toprule
Mode & Meaning and role in the thesis \\
\midrule
\Paxos{} & Original Multi-Paxos backend. It is the baseline because it is the replication path that \Mako{} previously used. \\
Multi-Raft & Direct Paxos-shaped \Raft{} replacement. Each partition lane has its own \Raft{} group, preserving the original parallel replication shape. \\
Single-Raft without \ReplayPool{} & Consolidated \Raft{} design that routes all partition commands through one process-wide \Raft{} log, with no follower replay offload. It exposes the bottleneck created when committed work feeds a narrow replay path. \\
Single-Raft with \ReplayPool{} & Consolidated \Raft{} design with the same process-wide \Raft{} log, but with follower replay moved into parallel replay workers. This is the engineered Single-Raft design evaluated as a viable replacement. \\
\bottomrule
\end{tabularx}
\end{table}

\Needspace{6\baselineskip}
The last two modes require the most care. Single-Raft without \ReplayPool{} is
not included because it is the final design. It is included because it isolates
the bottleneck that appears when all partition commands are consolidated into
one process-wide \Raft{} log and replay remains too narrow. Single-Raft with
\ReplayPool{} then tests the actual fix: keep \Raft{} apply lightweight and
move heavy follower replay into a parallel pool.

\section{What This Background Establishes}

The original \Mako{} path imposes three requirements on any Raft replacement
backend. First, the backend must preserve the transaction pipeline.
The leader should still be able to execute transactions, produce compact log
records, and submit those records for replication without turning consensus
into the transaction execution engine itself. \Raft{} should order records; it
should not force a redesign of how \Mako{} workers execute transactions.

Second, the backend must respect \Mako{}'s parallelism. The original
Multi-Paxos design was not only a correctness mechanism. It was also a
performance structure that allowed multiple worker lanes to produce and
replicate work in parallel. Multi-Raft preserves this structure directly by
using separate \Raft{} groups. Single-Raft changes the structure, so it must be
judged not only by the number of committed records at the leader, but also by
whether follower replay keeps up with those records.

Third, every Raft backend must fit the way \Mako{} submits work. In the
experiments used in this thesis, the benchmark drives transaction work through
the leader-side process. The original Paxos path did not expose a separate
dynamic leader-placement problem to Mako because the benchmark and replication
setup already used a fixed leader-side role. Raft does introduce that problem:
if \Raft{} elects a different process as leader, the consensus protocol may
still be correct, but the system no longer matches the leader-oriented
submission path being evaluated; local submissions at the benchmark process may
be refused or fail to represent the intended Paxos/Raft comparison. This is
the motivation for the preferred leader mechanism discussed later, and it
applies to both Multi-Raft and Single-Raft.

These requirements explain the structure of the remaining chapters. Chapter 3
shows how \Raft{} is placed behind the same replication boundary. Chapter 4
explains why leadership placement matters for \Mako{}'s leader-oriented
submission path. Chapter 5 explains why Single-Raft needs replay offload. The
evaluation then checks the same basic question from several angles: does the
\Raft{} backend preserve committed throughput, worker CPU use, and follower
replay progress well enough to remain in \Paxos{}'s performance class?
